\documentclass[12pt,a4paper]{article}

\usepackage[utf8]{inputenc}
\usepackage[T1]{fontenc}
\usepackage[polish]{babel}
\usepackage[a4paper,margin=2.5cm]{geometry}
\usepackage{enumitem}
\usepackage{hyperref}
\usepackage{titlesec}
\usepackage{graphicx}
\usepackage{tabularx}
\usepackage{booktabs}
\usepackage{listings}
\usepackage{xcolor}
\usepackage{float}
\usepackage{amsmath}
\usepackage{amssymb}
\usepackage{longtable}
\usepackage{multirow}

\definecolor{prologbg}{RGB}{245,245,245}
\definecolor{prologkw}{RGB}{0,0,180}
\definecolor{prologcomment}{RGB}{0,128,0}
\definecolor{prologstring}{RGB}{163,21,21}

\lstdefinestyle{prolog}{
    language=Prolog,
    backgroundcolor=\color{prologbg},
    basicstyle=\ttfamily\small,
    keywordstyle=\color{prologkw}\bfseries,
    commentstyle=\color{prologcomment}\itshape,
    stringstyle=\color{prologstring},
    breaklines=true,
    frame=single,
    numbers=left,
    numberstyle=\tiny\color{gray},
    tabsize=4,
    showstringspaces=false,
    captionpos=b,
    morekeywords={findall,forall,predsort,retractall,assertz,member,once,setup_call_cleanup,write,nl,read_line_to_string,atom_string,string_lower,normalize_space,ensure_loaded,use_module,dynamic,begin_tests,end_tests}
}

\newcommand{\projecttitle}{Podróżniczy System Ekspertowy -- Dokumentacja Techniczna}
\newcommand{\authorone}{Aleksandra Wieczorek}
\newcommand{\authortwo}{Wojciech Wójcik}
\newcommand{\projectdate}{Warszawa, 2026}

\setcounter{tocdepth}{2}

\begin{document}
\sloppy

\begin{titlepage}
    \centering

    \includegraphics[width=0.8\textwidth]{politechnika.png}\par
    \vspace{2.5cm}

    {\Large Dokumentacja techniczna projektu z~przedmiotu Systemy Ekspertowe\par}
    \vspace{1.2cm}

    {\Huge \bfseries \projecttitle \par}
    \vspace{2cm}

    {\Large Autorzy:\par}
    \vspace{0.5cm}
    {\large \authorone\par}
    \vspace{0.3cm}
    {\large \authortwo\par}

    \vfill

    {\large \projectdate\par}
\end{titlepage}

\tableofcontents
\newpage

\section{Wprowadzenie}

Niniejszy dokument stanowi dokumentację techniczną systemu ekspertowego wspomagającego wybór destynacji podróżniczej. System został zaimplementowany w~języku Prolog (SWI-Prolog) i~łączy w~sobie trzy techniki sztucznej inteligencji: klasyczne wnioskowanie regułowe, logikę rozmytą (fuzzy logic) oraz teorię zbiorów przybliżonych (rough sets) w~postaci algorytmu wyznaczania minimalnych reduktów.

System prowadzi interaktywną konsultację z~użytkownikiem, zbiera informacje o~preferencjach podróżniczych, a~następnie generuje ranking destynacji wraz z~uzasadnieniem rekomendacji oraz analizą minimalnych zbiorów atrybutów potrzebnych do podjęcia decyzji.

\subsection{Uruchomienie systemu}

System wymaga zainstalowanego środowiska SWI-Prolog. Uruchomienie odbywa się poleceniem:

\begin{lstlisting}[style=prolog, numbers=none]
swipl -q -s src/main.pl
\end{lstlisting}

Po załadowaniu modułów należy wywołać predykat \texttt{main.} w~konsoli Prologa. System zada użytkownikowi siedem pytań, a~następnie wyświetli wyniki.

Uruchomienie testów jednostkowych:

\begin{lstlisting}[style=prolog, numbers=none]
swipl -g "load_test_files([]), run_tests." -t halt tests/test_inference.pl
\end{lstlisting}

\section{Architektura systemu}

System posiada budowę modułową. Każdy moduł jest osobnym plikiem Prologa odpowiedzialnym za wyodrębniony obszar funkcjonalności. Główny plik \texttt{main.pl} ładuje wszystkie moduły i~definiuje przebieg konsultacji.

\subsection{Struktura katalogów}

\begin{verbatim}
sys-eks-podroze/
  src/
    main.pl           -- punkt wejścia, ładowanie modułów
    questions.pl      -- interfejs użytkownika, pytania
    destinations.pl   -- baza wiedzy o destynacjach
    rules.pl          -- reguły eksperckie z logiką rozmytą
    inference.pl      -- mechanizm wnioskowania i scoringu
    reducts.pl        -- algorytm minimalnych reduktów
  tests/
    test_inference.pl -- testy jednostkowe
  README.md           -- instrukcja uruchomienia
\end{verbatim}

\subsection{Diagram modułów}

Zależności między modułami przedstawiają się następująco:

\begin{verbatim}
  main.pl
    |--- destinations.pl   (baza faktów)
    |--- questions.pl       (interfejs użytkownika)
    |--- rules.pl           (reguły eksperckie)
    |       |--- korzysta z: destinations.pl, questions.pl
    |--- inference.pl       (mechanizm wnioskowania)
    |       |--- korzysta z: rules.pl, destinations.pl
    |--- reducts.pl         (redukty)
            |--- korzysta z: inference.pl, questions.pl
\end{verbatim}

\subsection{Przebieg konsultacji}

Predykat \texttt{main} wykonuje następujące kroki:

\begin{enumerate}
    \item \texttt{start\_consultation} -- czyszczenie poprzednich odpowiedzi, wyświetlenie informacji, zadanie 7~pytań.
    \item \texttt{show\_recommendations} -- obliczenie i~wyświetlenie rankingu destynacji.
    \item \texttt{show\_reducts} -- wyznaczenie i~wyświetlenie minimalnych reduktów (ścisłych).
    \item \texttt{show\_reduct\_recommendations} -- wyświetlenie rekomendacji przy użyciu każdego reduktu.
    \item \texttt{show\_reducts\_tolerant\_aggressive} -- redukty tolerancyjne (spadek do 25 pkt).
    \item \texttt{show\_used\_attributes} -- lista użytych atrybutów.
\end{enumerate}

\section{Baza wiedzy -- moduł \texttt{destinations.pl}}

Baza wiedzy zawiera fakty opisujące 13~destynacji podróżniczych oraz ich właściwości. Wiedza jest zorganizowana w~kilku kategoriach predykatów.

\subsection{Lista destynacji}

System zawiera następujące destynacje:

\begin{table}[H]
\centering
\begin{tabular}{ll}
\toprule
\textbf{Destynacja} & \textbf{Kraj} \\
\midrule
Barcelona & Hiszpania \\
Paryż & Francja \\
Rzym & Włochy \\
Zakopane & Polska \\
Praga & Czechy \\
Lizbona & Portugalia \\
Wiedeń & Austria \\
Budapeszt & Węgry \\
Madeira & Portugalia \\
Oslo & Norwegia \\
Santorini & Grecja \\
Kraków & Polska \\
Split & Chorwacja \\
\bottomrule
\end{tabular}
\caption{Destynacje w~bazie wiedzy}
\end{table}

\subsection{Kategorie atrybutów destynacji}

Każda destynacja jest opisana za pomocą następujących predykatów:

\begin{itemize}
    \item \texttt{destination(D)} -- deklaracja destynacji,
    \item \texttt{country(D, Kraj)} -- kraj, w~którym znajduje się destynacja,
    \item \texttt{feature(D, Cecha)} -- cechy binarne (np. \texttt{ma\_plaze}, \texttt{ma\_zabytki}, \texttt{dobra\_komunikacja}, \texttt{wysoka\_temperatura\_latem}, \texttt{kultura}, \texttt{natura}, \texttt{aktywny\_wypoczynek}, \texttt{bezpieczny}, \texttt{tani}, \texttt{dobry\_na\_weekend}, \texttt{odpowiedni\_dla\_rodzin}, \texttt{romantyczny}, \texttt{rozrywka}, \texttt{chlodny\_klimat}),
    \item \texttt{cost(D, Koszt)} -- poziom kosztów: \texttt{niski}, \texttt{sredni}, \texttt{wysoki},
    \item \texttt{climate(D, Klimat)} -- klimat: \texttt{cieply}, \texttt{umiarkowany}, \texttt{chlodny},
    \item \texttt{good\_length(D, Dlugosc)} -- optymalna długość wyjazdu,
    \item \texttt{transport\_ok(D, Transport)} -- dostępne środki transportu,
    \item \texttt{good\_for(D, Towarzystwo)} -- odpowiedni typ towarzystwa,
    \item \texttt{activity\_level(D, Poziom)} -- poziom aktywności: \texttt{niski}, \texttt{sredni}, \texttt{wysoki}.
\end{itemize}

\subsection{Przykład reprezentacji destynacji}

\begin{lstlisting}[style=prolog, caption={Reprezentacja Barcelony w~bazie wiedzy}]
destination(barcelona).
country(barcelona, hiszpania).

feature(barcelona, ma_plaze).
feature(barcelona, ma_zabytki).
feature(barcelona, dobra_komunikacja).
feature(barcelona, wysoka_temperatura_latem).
feature(barcelona, rozrywka).

cost(barcelona, sredni).
climate(barcelona, cieply).

good_length(barcelona, kilka_dni).
good_length(barcelona, tydzien).

transport_ok(barcelona, samolot).

good_for(barcelona, znajomi).
good_for(barcelona, para).

activity_level(barcelona, sredni).
\end{lstlisting}

\section{Interfejs użytkownika -- moduł \texttt{questions.pl}}

Moduł \texttt{questions.pl} odpowiada za prowadzenie dialogu z~użytkownikiem. Odpowiedzi użytkownika są przechowywane jako dynamiczne fakty \texttt{answer(Atrybut, Wartość)}.

\subsection{Zbierane atrybuty}

System zadaje użytkownikowi 7~pytań. Każde pytanie ma zdefiniowany zbiór dopuszczalnych odpowiedzi:

\begin{table}[H]
\centering
\small
\begin{tabularx}{\textwidth}{lX}
\toprule
\textbf{Atrybut} & \textbf{Dopuszczalne wartości} \\
\midrule
\texttt{typ\_wypoczynku} & plażowy, miejski, górski, aktywny, kulturowy, relaksacyjny \\
\texttt{budzet} & bardzo\_niski, niski, średni, wysoki, bardzo\_wysoki \\
\texttt{klimat} & bardzo\_ciepły, ciepły, umiarkowany, chłodny, bardzo\_chłodny \\
\texttt{srodek\_transportu} & samolot, samochód, pociąg, dowolny \\
\texttt{dlugosc\_wyjazdu} & krótki, weekend, kilka\_dni, tydzień, długi\_pobyt \\
\texttt{towarzystwo} & solo, para, rodzina, znajomi \\
\texttt{poziom\_aktywnosci} & bardzo\_niski, niski, średni, wysoki, bardzo\_wysoki \\
\bottomrule
\end{tabularx}
\caption{Atrybuty preferencji użytkownika i~ich dopuszczalne wartości}
\end{table}

\subsection{Mechanizm walidacji}

Odpowiedzi użytkownika są normalizowane (konwersja do małych liter, usunięcie nadmiarowych spacji) za pomocą predykatu \texttt{normalize\_input}. Jeśli użytkownik poda niepoprawną odpowiedź, system wyświetla komunikat o~błędzie i~ponawia pytanie.

\begin{lstlisting}[style=prolog, caption={Normalizacja odpowiedzi użytkownika}]
normalize_input(Input, Atom) :-
    string_lower(Input, Lower),
    normalize_space(string(Clean), Lower),
    atom_string(Atom, Clean).
\end{lstlisting}

\section{Reguły eksperckie i~logika rozmyta -- moduł \texttt{rules.pl}}

Moduł \texttt{rules.pl} stanowi serce wiedzy eksperckiej systemu. Reguły łączą preferencje użytkownika z~właściwościami destynacji, wykorzystując logikę rozmytą do obsługi nieostrych pojęć.

\subsection{Struktura reguł}

Wszystkie reguły eksperckie mają postać predykatu:

\begin{lstlisting}[style=prolog, numbers=none]
preference_point(Destynacja, Punkty, Powod).
\end{lstlisting}

Każda reguła, która jest spełniona dla danej destynacji, przyznaje określoną liczbę punktów wraz z~tekstowym uzasadnieniem. Punkty ze wszystkich spełnionych reguł są sumowane przez mechanizm wnioskowania.

\subsection{Logika rozmyta -- funkcje podobieństwa}

System wykorzystuje \textbf{funkcje podobieństwa} (similarity functions), które przypisują stopień przynależności $\mu \in [0, 1]$ do każdego dopasowania między preferencją użytkownika a~cechą destynacji. Obliczenie punktów rozmytych realizowane jest wzorem:

\begin{equation}
\text{Points} = \text{round}(\text{MaxPoints} \times \mu)
\end{equation}

gdzie $\mu$ to stopień podobieństwa, a~$\text{MaxPoints}$ to maksymalna liczba punktów za daną kategorię.

\begin{lstlisting}[style=prolog, caption={Funkcja obliczania punktów rozmytych}]
fuzzy_points(MaxPoints, Degree, Points) :-
    PointsFloat is MaxPoints * Degree,
    Points is round(PointsFloat).
\end{lstlisting}

W~systemie zdefiniowano 5~funkcji podobieństwa:

\subsubsection{Podobieństwo typu wypoczynku (\texttt{type\_similarity})}

Mapuje preferencje wypoczynku na cechy destynacji. Stopień przynależności zależy od tego, jak dobrze destynacja pasuje do wybranego typu:

\begin{table}[H]
\centering
\small
\begin{tabular}{llc}
\toprule
\textbf{Preferencja} & \textbf{Warunek} & $\mu$ \\
\midrule
plażowy & destynacja ma plaże & 1.0 \\
plażowy & wysoka temp., brak plaży & 0.7 \\
miejski & komunikacja + zabytki & 1.0 \\
miejski & tylko komunikacja lub zabytki & 0.7 \\
górski & destynacja ma góry & 1.0 \\
górski & natura, brak gór & 0.7 \\
aktywny & aktywny wypoczynek & 1.0 \\
aktywny & góry, brak aktywnego wyp. & 0.7 \\
aktywny & natura (brak gór i~aktywnego) & 0.5 \\
kulturowy & destynacja ma kulturę & 1.0 \\
kulturowy & zabytki, brak kultury & 0.8 \\
relaksacyjny & plaże lub natura & 1.0 \\
relaksacyjny & bezpieczny (brak plaży/natury) & 0.7 \\
\bottomrule
\end{tabular}
\caption{Stopnie przynależności dla typu wypoczynku (maks.\ 20 pkt)}
\end{table}

\subsubsection{Podobieństwo budżetu (\texttt{budget\_similarity})}

Macierz podobieństwa między preferowanym budżetem użytkownika a~kosztem destynacji:

\begin{table}[H]
\centering
\begin{tabular}{l|ccc}
\toprule
\textbf{Preferencja / Koszt} & \textbf{niski} & \textbf{średni} & \textbf{wysoki} \\
\midrule
bardzo\_niski & 1.0 & 0.6 & 0.1 \\
niski & 1.0 & 0.75 & 0.2 \\
średni & 0.4 & 1.0 & 0.45 \\
wysoki & 0.1 & 0.65 & 1.0 \\
bardzo\_wysoki & -- & 0.4 & 1.0 \\
\bottomrule
\end{tabular}
\caption{Macierz podobieństwa budżetu (maks.\ 15 pkt)}
\end{table}

\subsubsection{Podobieństwo klimatu (\texttt{climate\_similarity})}

\begin{table}[H]
\centering
\begin{tabular}{l|ccc}
\toprule
\textbf{Preferencja / Klimat} & \textbf{ciepły} & \textbf{umiarkowany} & \textbf{chłodny} \\
\midrule
bardzo\_ciepły & 1.0 & 0.6 & 0.1 \\
ciepły & 1.0 & 0.7 & 0.2 \\
umiarkowany & 0.6 & 1.0 & 0.6 \\
chłodny & 0.1 & 0.7 & 1.0 \\
bardzo\_chłodny & 0.1 & 0.5 & 1.0 \\
\bottomrule
\end{tabular}
\caption{Macierz podobieństwa klimatu (maks.\ 15 pkt)}
\end{table}

\subsubsection{Podobieństwo długości wyjazdu (\texttt{length\_similarity})}

\begin{table}[H]
\centering
\begin{tabular}{l|cccc}
\toprule
\textbf{Pref. / Długość} & \textbf{weekend} & \textbf{kilka\_dni} & \textbf{tydzień} & \textbf{długi\_pobyt} \\
\midrule
krótki & 1.0 & 0.6 & 0.1 & -- \\
weekend & 1.0 & 0.8 & 0.2 & -- \\
kilka\_dni & 0.7 & 1.0 & 0.7 & -- \\
tydzień & -- & 0.6 & 1.0 & 0.7 \\
długi\_pobyt & -- & -- & 0.6 & 1.0 \\
\bottomrule
\end{tabular}
\caption{Macierz podobieństwa długości wyjazdu (maks.\ 15 pkt)}
\end{table}

\subsubsection{Podobieństwo poziomu aktywności (\texttt{activity\_similarity})}

\begin{table}[H]
\centering
\begin{tabular}{l|ccc}
\toprule
\textbf{Preferencja / Poziom} & \textbf{niski} & \textbf{średni} & \textbf{wysoki} \\
\midrule
bardzo\_niski & 1.0 & 0.5 & 0.1 \\
niski & 1.0 & 0.75 & 0.2 \\
średni & 0.4 & 1.0 & 0.5 \\
wysoki & 0.1 & 0.6 & 1.0 \\
bardzo\_wysoki & -- & 0.4 & 1.0 \\
\bottomrule
\end{tabular}
\caption{Macierz podobieństwa poziomu aktywności (maks.\ 10 pkt)}
\end{table}

\subsection{Rozkład punktów w~regułach}

Reguły przyznają punkty w~następujących kategoriach:

\begin{table}[H]
\centering
\begin{tabular}{lcc}
\toprule
\textbf{Kategoria reguły} & \textbf{Maks.\ punkty} & \textbf{Typ} \\
\midrule
Typ wypoczynku & 20 & rozmyta \\
Budżet & 15 & rozmyta \\
Klimat & 15 & rozmyta \\
Długość wyjazdu & 15 & rozmyta \\
Transport & 10 & binarna \\
Towarzystwo & 10 & binarna \\
Poziom aktywności & 10 & rozmyta \\
\midrule
\textit{Reguły specjalistyczne (każda)} & \textit{10} & \textit{binarna} \\
\bottomrule
\end{tabular}
\caption{Maksymalna liczba punktów w~poszczególnych kategoriach reguł}
\end{table}

\subsection{Reguły specjalistyczne}

Oprócz reguł opartych na funkcjach podobieństwa system zawiera 5~reguł specjalistycznych, które przyznają dodatkowe punkty bonusowe (po 10 pkt) w~szczególnych sytuacjach:

\begin{enumerate}
    \item \textbf{Tani wyjazd} -- jeśli użytkownik preferuje niski budżet i~destynacja ma cechę \texttt{tani}.
    \item \textbf{Weekend} -- jeśli użytkownik planuje wyjazd weekendowy i~destynacja ma cechę \texttt{dobry\_na\_weekend}.
    \item \textbf{Rodzina} -- jeśli użytkownik podróżuje z~rodziną i~destynacja jest \texttt{odpowiedni\_dla\_rodzin}.
    \item \textbf{Podróż solo} -- jeśli użytkownik podróżuje sam i~destynacja jest \texttt{bezpieczny}.
    \item \textbf{Ciepła pogoda} -- jeśli użytkownik preferuje ciepły klimat i~destynacja ma \texttt{wysoka\_temperatura\_latem}.
\end{enumerate}

\begin{lstlisting}[style=prolog, caption={Przykłady reguł specjalistycznych}]
preference_point(D, 10, 'dobra opcja dla taniego wyjazdu') :-
    answer(budzet, niski),
    feature(D, tani).

preference_point(D, 10, 'dobra opcja na weekend') :-
    answer(dlugosc_wyjazdu, weekend),
    feature(D, dobry_na_weekend).

preference_point(D, 10, 'bezpieczna opcja dla podrozy solo') :-
    answer(towarzystwo, solo),
    feature(D, bezpieczny).
\end{lstlisting}

\section{Mechanizm wnioskowania -- moduł \texttt{inference.pl}}

Moduł \texttt{inference.pl} implementuje silnik inferencyjny systemu, który agreguje wyniki z~reguł eksperckich i~generuje ranking destynacji.

\subsection{Algorytm wnioskowania}

Dla każdej destynacji $D$ w~bazie wiedzy system wykonuje następujące kroki:

\begin{enumerate}
    \item Za pomocą \texttt{findall} zbiera wszystkie pary $(\text{Punkty}, \text{Powód})$ z~reguł \texttt{preference\_point(D, Punkty, Powód)}, które są spełnione przy aktualnych odpowiedziach użytkownika.
    \item Sumuje punkty ze wszystkich spełnionych reguł: $\text{Score}(D) = \sum_{i} \text{Points}_i$.
    \item Ekstrahuje listę powodów (uzasadnień) z~dopasowanych reguł.
    \item Sortuje destynacje malejąco według wyniku.
\end{enumerate}

\begin{lstlisting}[style=prolog, caption={Główny predykat rekomendacji}]
recommendation(Destination, Score, Reasons) :-
    destination(Destination),
    findall(
        Points-Reason,
        preference_point(Destination, Points, Reason),
        Matches
    ),
    Matches \= [],
    sum_points(Matches, Score),
    extract_reasons(Matches, Reasons).

all_recommendations(Sorted) :-
    findall(
        Score-Destination-Reasons,
        recommendation(Destination, Score, Reasons),
        Recommendations
    ),
    predsort(compare_recommendations, Recommendations, Sorted).
\end{lstlisting}

\subsection{Sortowanie wyników}

Destynacje są sortowane za pomocą \texttt{predsort} z~dedykowanym komparatorem, który porównuje wyniki malejąco (destynacja z~najwyższym wynikiem jest pierwsza):

\begin{lstlisting}[style=prolog, caption={Komparator do sortowania rekomendacji}]
compare_recommendations(Order, Score1-_-_, Score2-_-_) :-
    (   Score1 > Score2 -> Order = <
    ;   Score1 < Score2 -> Order = >
    ;   Order = =
    ).
\end{lstlisting}

\subsection{Eksplanacja wyników}

System dostarcza eksplanację (wyjaśnienie) każdej rekomendacji w~postaci listy powodów. Każdy powód to tekstowy opis reguły, która została spełniona. Dzięki temu użytkownik rozumie, dlaczego dana destynacja znalazła się na określonej pozycji w~rankingu.

Przykładowy wynik:
\begin{verbatim}
1. barcelona - dopasowanie: 114 pkt
   Powody:
   - pasuje do wypoczynku plazowego
   - pasuje do wybranego budzetu
   - pasuje do preferowanego klimatu
   - pasuje do dlugosci wyjazdu
   - pasuje do preferowanego transportu
   - pasuje do typu towarzystwa
   - pasuje do poziomu aktywnosci
   - ciepla destynacja z dobra pogoda latem
\end{verbatim}

\section{Redukty -- moduł \texttt{reducts.pl}}

Moduł \texttt{reducts.pl} implementuje algorytm wyznaczania \textbf{minimalnych reduktów} oparty na teorii zbiorów przybliżonych (rough set theory). Redukt to minimalny podzbiór odpowiedzi użytkownika, który zachowuje tę~samą najlepszą rekomendację co pełny zbiór atrybutów.

\subsection{Definicja formalna}

Niech $A = \{a_1, a_2, \ldots, a_n\}$ będzie zbiorem odpowiedzi użytkownika, $R(A)$ będzie najlepszą rekomendacją przy użyciu pełnego zbioru $A$, a~$S(A)$ -- wynikiem punktowym tej rekomendacji.

Podzbiór $P \subseteq A$ jest \textbf{reduktem}, jeśli:
\begin{enumerate}
    \item $R(P) = R(A)$ oraz $S(P) = S(A)$ (zachowanie rekomendacji i~wyniku),
    \item $\forall P' \subset P,\; P' \neq \emptyset: R(P') \neq R(A) \lor S(P') \neq S(A)$ (minimalność -- żaden właściwy podzbiór nie spełnia warunku 1).
\end{enumerate}

\subsection{Algorytm}

\begin{lstlisting}[style=prolog, caption={Wyznaczanie minimalnych reduktów}]
minimal_reduct_for(Destination, Score, Answers, Reduct) :-
    subset_of_answers(Reduct, Answers),
    Reduct \= [],
    with_answers(Reduct,
        best_recommendation(Destination, Score, _)),
    \+ (
        proper_subset_of_answers(ProperSubset, Reduct),
        ProperSubset \= [],
        with_answers(ProperSubset,
            best_recommendation(Destination, Score, _))
    ).
\end{lstlisting}

Algorytm działa następująco:
\begin{enumerate}
    \item Generuje wszystkie niepuste podzbiory odpowiedzi użytkownika.
    \item Dla każdego podzbioru tymczasowo zastępuje odpowiedzi użytkownika podzbiorem (predykat \texttt{with\_answers}).
    \item Sprawdza, czy najlepsza rekomendacja jest taka sama jak przy pełnym zbiorze odpowiedzi.
    \item Weryfikuje minimalność -- upewnia się, że żaden właściwy podzbiór danego podzbioru nie daje takiego samego wyniku.
\end{enumerate}

\subsection{Bezpieczna podmiana odpowiedzi}

Predykat \texttt{with\_answers} tymczasowo podmienia odpowiedzi użytkownika w~bazie dynamicznej, wykonuje zadany cel, a~następnie przywraca oryginalne odpowiedzi. Wykorzystuje \texttt{setup\_call\_cleanup} dla bezpieczeństwa:

\begin{lstlisting}[style=prolog, caption={Bezpieczna podmiana kontekstu odpowiedzi}]
with_answers(Answers, Goal) :-
    answered_answers(OriginalAnswers),
    setup_call_cleanup(
        (retractall(answer(_, _)),
         assert_answer_terms(Answers)),
        Goal,
        (retractall(answer(_, _)),
         assert_answer_terms(OriginalAnswers))
    ).
\end{lstlisting}

\subsection{Redukty tolerancyjne}

Oprócz reduktów ścisłych system oferuje \textbf{redukty tolerancyjne}, które akceptują podzbiory atrybutów nawet jeśli wynik punktowy nieznacznie spadnie. Parametr \texttt{MaxDrop} określa maksymalny dopuszczalny spadek wyniku.

Podzbiór $P$ jest reduktem tolerancyjnym z~tolerancją $\delta$, jeśli:
\begin{equation}
R(P) = R(A) \quad \lor \quad S(P) \geq S(A) - \delta
\end{equation}

System domyślnie prezentuje redukty z~tolerancją $\delta = 25$ punktów (\texttt{show\_reducts\_tolerant\_aggressive}), co pozwala na bardziej agresywną redukcję atrybutów.

\subsection{Złożoność obliczeniowa}

Algorytm wymaga wygenerowania wszystkich podzbiorów zbioru odpowiedzi, co daje złożoność $O(2^n)$, gdzie $n$ to liczba atrybutów. Dla $n = 7$ (128 podzbiorów) jest to obliczeniowo wykonalne. Dla każdego podzbioru wykonywane jest wnioskowanie, co daje całkowitą złożoność:

\begin{equation}
O(2^n \cdot m \cdot r)
\end{equation}

gdzie $m$ to liczba destynacji (13), a~$r$ to liczba reguł.

\section{Scenariusze użycia}

Poniżej przedstawiono cztery przykładowe scenariusze konsultacji demonstrujące różne aspekty działania systemu.

\subsection{Scenariusz 1: Wyjazd plażowy (para, samolot)}

\begin{table}[H]
\centering
\begin{tabular}{ll}
\toprule
\textbf{Atrybut} & \textbf{Wartość} \\
\midrule
Typ wypoczynku & plażowy \\
Budżet & średni \\
Klimat & ciepły \\
Transport & samolot \\
Długość & tydzień \\
Towarzystwo & para \\
Aktywność & średni \\
\bottomrule
\end{tabular}
\end{table}

\textbf{Oczekiwany wynik:} Barcelona na pierwszym miejscu ($\sim$114 pkt) -- pełne dopasowanie plaż, klimatu, budżetu i~długości. Lizbona i~Split na dalszych pozycjach. Redukt agresywny redukuje wynik do $\sim$89 pkt przy zachowaniu rekomendacji.

\subsection{Scenariusz 2: Weekend miejski (niski budżet, pociąg)}

\begin{table}[H]
\centering
\begin{tabular}{ll}
\toprule
\textbf{Atrybut} & \textbf{Wartość} \\
\midrule
Typ wypoczynku & miejski \\
Budżet & niski \\
Klimat & umiarkowany \\
Transport & pociąg \\
Długość & weekend \\
Towarzystwo & para \\
Aktywność & niski \\
\bottomrule
\end{tabular}
\end{table}

\textbf{Oczekiwany wynik:} Praga na pierwszym miejscu ($\sim$127 pkt) -- tania destynacja z~zabytkami, dobra komunikacja, dostępna pociągiem. Budapeszt na drugiej pozycji ($\sim$115 pkt). Redukt pokazuje, że mniej atrybutów wystarczy do identyfikacji Pragi.

\subsection{Scenariusz 3: Górski aktywny (rodzina, samochód)}

\begin{table}[H]
\centering
\begin{tabular}{ll}
\toprule
\textbf{Atrybut} & \textbf{Wartość} \\
\midrule
Typ wypoczynku & aktywny \\
Budżet & średni \\
Klimat & umiarkowany \\
Transport & samochód \\
Długość & kilka\_dni \\
Towarzystwo & rodzina \\
Aktywność & wysoki \\
\bottomrule
\end{tabular}
\end{table}

\textbf{Oczekiwany wynik:} Zakopane na pierwszym miejscu ($\sim$116 pkt) -- góry, aktywny wypoczynek, odpowiednie dla rodzin, dostępne samochodem.

\subsection{Scenariusz 4: Luksusowa plaża (para, minimalny zestaw atrybutów)}

\begin{table}[H]
\centering
\begin{tabular}{ll}
\toprule
\textbf{Atrybut} & \textbf{Wartość} \\
\midrule
Typ wypoczynku & plażowy \\
Budżet & wysoki \\
Towarzystwo & para \\
\bottomrule
\end{tabular}
\end{table}

\textbf{Oczekiwany wynik:} Santorini ($\sim$45 pkt) -- romantyczna wyspa z~plażami, wysoki budżet dopasowany. Redukt agresywny pokazuje, że wystarczają zaledwie 3~atrybuty ($\sim$20 pkt) do wskazania tej destynacji.

\section{Testy jednostkowe}

System posiada zestaw testów jednostkowych zdefiniowanych w~pliku \texttt{tests/test\_inference.pl}, wykorzystujących wbudowany framework testowy SWI-Prolog.

\subsection{Przypadki testowe}

\begin{table}[H]
\centering
\small
\begin{tabularx}{\textwidth}{lX}
\toprule
\textbf{Nazwa testu} & \textbf{Opis} \\
\midrule
\texttt{beach\_recommendation\_exists} & Sprawdza, czy dla scenariusza plażowego generowane są rekomendacje (lista niepusta). \\
\texttt{barcelona\_should\_score\_for\_beach\_trip} & Weryfikuje, że Barcelona uzyskuje wynik $> 0$ dla wyjazdu plażowego. \\
\texttt{fuzzy\_climate\_match\_scores\_partially} & Testuje, że Praga przy preferencji klimatu ciepłego uzyskuje dokładnie 11 pkt (częściowe dopasowanie rozmyte: $\text{round}(15 \times 0.7) = 11$). \\
\texttt{minimal\_reducts\_preserve\_best} & Sprawdza, że redukt zachowuje tę~samą rekomendację i~wynik co pełny zbiór atrybutów. \\
\texttt{used\_attributes\_are\_detected} & Weryfikuje poprawne wykrywanie użytych atrybutów. \\
\bottomrule
\end{tabularx}
\caption{Przypadki testowe systemu}
\end{table}

\subsection{Przykład testu rozmytego}

Test \texttt{fuzzy\_climate\_match\_scores\_partially} ilustruje działanie logiki rozmytej:

\begin{lstlisting}[style=prolog, caption={Test częściowego dopasowania klimatu}]
test(fuzzy_climate_match_scores_partially) :-
    retractall(answer(_, _)),
    assertz(answer(klimat, cieply)),
    recommendation(praga, Score, _),
    Score =:= 11.
\end{lstlisting}

Praga ma klimat \texttt{umiarkowany}. Przy preferencji \texttt{cieply} stopień podobieństwa wynosi $\mu = 0.7$, co daje $\text{round}(15 \times 0.7) = 11$ punktów.

\section{Podsumowanie}

Zaimplementowany system ekspertowy łączy trzy techniki sztucznej inteligencji w~spójnym narzędziu rekomendacyjnym:

\begin{enumerate}
    \item \textbf{Klasyczne wnioskowanie regułowe} -- reguły produkcji \texttt{preference\_point} łączą preferencje użytkownika z~cechami destynacji, umożliwiając agregację punktów i~generowanie rankingu.

    \item \textbf{Logika rozmyta} -- funkcje podobieństwa (\texttt{type\_similarity}, \texttt{budget\_similarity}, \texttt{climate\_similarity}, \texttt{length\_similarity}, \texttt{activity\_similarity}) przyznają częściowe punkty za niepełne dopasowania, co pozwala na płynną ocenę zamiast binarnej klasyfikacji.

    \item \textbf{Teoria zbiorów przybliżonych} -- algorytm minimalnych reduktów identyfikuje najmniejsze podzbiory atrybutów wystarczające do zachowania rekomendacji, co demonstruje, które preferencje użytkownika są rzeczywiście decydujące.
\end{enumerate}

System zawiera 13~destynacji, 7~atrybutów preferencji z~5~rozmytymi funkcjami podobieństwa, reguły specjalistyczne oraz dwa tryby redukcji atrybutów (ścisły i~tolerancyjny). Modułowa architektura ułatwia rozszerzanie bazy wiedzy o~nowe destynacje i~reguły bez modyfikacji mechanizmu wnioskowania.

\end{document}
